\section{Rigidity-bridge certificates and mainline checklist (audit)}
\label{app:rigidity_bridge_checklist}

\AuditTag This appendix records a compact checklist for ``rigidity bridges'': auditable certificate forms that turn a narrative step into a verifiable implication with an explicit minimal input set.
It is reader-facing audit infrastructure and is not used as a premise in theorem-level folding proofs.

\subsection{Certificate forms (RB-A/B/C/D)}
\label{subsec:rb_certificate_forms}

We use four certificate templates:
\begin{itemize}
  \item \textbf{RB-A (finite-family minimization).}
  A finite candidate family $\mathcal{F}$, an explicit objective functional $J:\mathcal{F}\to\RR$, and a deterministic tie-break rule; the output is a unique (or near-unique with a gap) minimizer $f_\star=\arg\min_{\mathcal{F}}J$.
  \item \textbf{RB-B (incompatibility / obstruction).}
  A ``bad'' structure would force an obstruction (e.g.\ an interior pole) that is incompatible with an established analytic domain or stability certificate; hence the bad structure is excluded.
  \item \textbf{RB-C (counting/classification rigidity).}
  A finite classification, counting identity, or image--preimage structure forces a unique decomposition (e.g.\ $64\to 21$, $18\oplus 3$).
  \item \textbf{RB-D (gap-stability / robustness).}
  A quantitative gap or sensitivity bound shows that the conclusion persists under bounded counterfactual families or perturbations, upgrading ``accidental alignment'' to a robust closure.
\end{itemize}

\subsection{Mainline checklist (where each bridge is realized in this paper)}
\label{subsec:rb_mainline_checklist}

\begin{center}
\scriptsize
\setlength{\tabcolsep}{6pt}
\renewcommand{\arraystretch}{1.2}
\begin{tabular}{p{0.18\textwidth} p{0.18\textwidth} p{0.14\textwidth} p{0.18\textwidth} p{0.24\textwidth}}
\toprule
mainline jump & closure output & RB form & minimal inputs & where in this paper \\
\midrule
Tick $\to$ CAP &
deterministic finite-family closure rule &
RB-A &
tick + bounded complexity + tie-break &
Axioms~\ref{ax:readout_sequentiality}, \ref{ax:cap}; Appendix~\ref{app:tick_cap_derivation}; Appendix~\ref{app:cap_audit_template} \\

CAP $\to$ golden branch &
finite-depth least-discrepancy selection of $\alpha=\varphi^{-1}$ &
RB-A/RB-D &
continued-fraction proxy + audited discrepancy bound &
Proposition~\ref{prop:golden_least_discrepancy}; Subsubsection~\ref{subsubsec:discrepancy_certificates}; Appendix~\ref{app:discrepancy_ostrowski} \\

golden branch $\to$ $\varphi$-grammar &
admissible set $X_m$ with $|X_m|=F_{m+2}$ &
RB-C &
Zeckendorf admissibility &
Lemma~\ref{lem:xm_fib}; Appendix~\ref{app:protocol_primitives} \\

$\varphi$-grammar $\to$ $\pi$-closure &
cyclic/boundary split and $18\oplus 3$ at $m=6$ &
RB-C &
wrap-around constraint + Fibonacci counts &
Proposition~\ref{prop:cyc_bdry_size}; Corollary~\ref{prop:cyc_bdry_6} \\

$(\varphi,\pi)\to \mathrm{e}$ stability &
zeta/Abel normalization and pole-barrier template &
RB-B/RB-D &
unit-disk holomorphy + Abel path conventions &
Section~\ref{subsec:e_channel}; Appendix~\ref{app:protocol_primitives}; Appendix~\ref{app:abel_finite_part_notes}; Appendix~\ref{app:trace_pole_barrier_template} \\

$\pi$-closure $\to$ anchor &
minimal explicit screen anchor $(m,n)=(6,3)$ and audited addressing choice &
RB-A/RB-C &
finite screen family + cardinality match + tie-break &
Section~\ref{sec:folding_core}; Remark~\ref{rem:balanced_coupling_convention}; Table~\ref{tab:addressing_selection} \\

anchor $\to$ gauge/holonomy &
finite connection/holonomy invariants modulo local relabeling &
RB-C/RB-D &
local fiber relabeling + loop invariants + bounded families &
Section~\ref{sec:protocol_connections_holonomy}; Appendix~\ref{app:holonomy_sweeps_extended} \\

gauge $\to$ SM labeling/mass &
bounded-family labeling closures at the anchor &
RB-A/RB-D &
explicit candidate families + gap/robustness audits &
Section~\ref{sec:sm_labeling_closure}; Appendix~\ref{app:generated_tables}; Appendix~\ref{app:closure_audit_details} \\

overhead $\to$ gravity/dynamics &
overhead-to-lapse/potential closure with error control &
RB-B/RB-D &
error budget + counterfactual regularization sweeps &
Appendix~\ref{app:overhead_to_gravity_closure}; Appendix~\ref{app:protocol_to_continuum_error_control} \\
\bottomrule
\end{tabular}
\end{center}

\subsection{A minimal workflow for upgrading intuition to audit statements}
\label{subsec:rb_workflow}

\AuditTag When adding an interface intuition, we recommend the following minimal workflow:
\begin{itemize}
  \item \textbf{Step 0 (decide the layer).} Decide whether the target sentence belongs to the theorem layer (\MathTag), the operational dictionary (\InterfaceTag), external calibration (\MatchTag), or audit/provenance (\AuditTag).
  \item \textbf{Step 1 (objectify and constrain).} Replace narrative words (``stable'', ``unique'', ``memoryless'') by explicit objects and constraints (functional equation, finite minimization, counting identity, or obstruction).
  \item \textbf{Step 2 (choose a certificate form).} Record which RB form (A/B/C/D) certifies the step, and state the minimal inputs used.
  \item \textbf{Step 3 (separate matching inputs).} Move unit choices, reference scales, and external targets into \MatchTag and/or the inference ledger, never as premises.
  \item \textbf{Step 4 (audit hooks).} Record the candidate family, objective, tie-break, and reproduction script entry points when applicable.
\end{itemize}


